\section{Related Work}
\label{sec:related}
Every ingredient \sys{} uses has a precedent: signals that trail data, in-network
multicast and reduction, and device-side barriers are all established. What we claim is
new is their combination and placement for cross-node EP synchronization, so this
section is organized around drawing that boundary precisely rather than around
categories of technique.

\paragraph{EP communication and its data path.}
DeepEP~\cite{deepep}, NCCL EP~\cite{ncclep}, and GPU-initiated
networking~\cite{gin} supply the setting, the baselines, and the transport primitives
we measure. We do not replace their data path; we relocate the coordination around it.
SwiftEP~\cite{swiftep} reduces copies and GPU overhead through buffer fusion and TMA
offloading, and COMET~\cite{comet} overlaps computation with communication at fine
grain. These optimize how the payload moves and when compute proceeds against it; our
question is when a rank may begin uploading and when a destination may be told its
input is complete. The two are composable, but improvements to byte movement cannot be
counted as control-traversal savings, and we keep the arithmetic and transferred bytes
fixed across any comparison.

\paragraph{Ordering inside an already-started transfer.}
Two recent systems are close enough to require careful separation, and in both cases
the distinction is the scope of what gets reordered.
Perseus~\cite{perseus} decouples each tile's PUT from its completion signal and uses a
NIC fence flag to preserve order, removing serialization in the proxy and NIC pipeline.
That serialization lives \emph{inside} an RDMA transfer that has already been
initiated; Perseus does not aggregate readiness across ranks, does not buffer early
data in the fabric, and does not generate per-destination completion.
UBEP~\cite{ubep} decomposes a BSP Dispatch pipeline on a production superpod,
substituting point-to-point signals and Data-as-Flag for its \texttt{SyncAll} and
explicit barriers. This too addresses stage synchronization within the EP data pipeline,
and it depends on pre-partitioned receive addresses and 512\,B atomic load/store over
its scale-up fabric. Both works corroborate that these costs are worth attacking. What
neither provides is readiness aggregated in the network, admission of early data against
bounded staging, or a fabric-side completion notification per destination.

\paragraph{In-network computation for MoE.}
DySHARP~\cite{dysharp} accelerates MoE with dynamic in-switch addressing, multicast and
reduction, and token-centric fusion, and MultiWrite~\cite{multiwrite} eliminates
duplicate upstream traffic in many-to-many communication via multicast write semantics.
These give us a data plane we can reuse, and we deliberately claim none of their
bandwidth benefit: our contribution concerns the synchronization critical path
surrounding those actions, not the bytes they save. Neither performs dynamic readiness
aggregation, staged release, or destination completion determination.

\paragraph{Collective offload and synchronization primitives.}
SHARP~\cite{sharp} offloads reduction and broadcast to the network, and
SwitchML~\cite{switchml} co-designs in-switch gradient aggregation with endpoint
protocols; both target collectives whose participant set and traffic volume are
established in advance, whereas MoE presents a readiness set and per-destination traffic
plan that change with routing. EPIC~\cite{epic} offers programmable protocol and
resource abstractions for Ethernet in-network computing and could plausibly host our
mechanism; our contribution is the MoE-specific event semantics and timing rather than a
general INC architecture. GPU-initiated networking~\cite{gin} exposes PUT, signal, and
fence primitives to the device, which is precisely the substrate the baseline tail is
built from, but it does not itself aggregate readiness across ranks or evaluate a
dynamic per-destination completion condition. PFC~\cite{pfc} can protect the staging
path at link level while supplying neither generation semantics nor an end-to-end
visibility contract.

\paragraph{Delay decomposition and observability.}
Swift~\cite{swift} decomposes an observed round trip into endpoint and fabric components
for congestion control; we borrow the decomposition discipline and propose no congestion
control algorithm. FabricPerf~\cite{fabricperf} develops fine-grained observation of GPU
scale-up communication, which informs our event-based methodology, though the
\ready{}/\data{}/\done{} events on cross-node Direct required separate instrumentation.
